\emph{Teorema kejangkeping} minangka kosok baliné \emph{teorema
kesahihan}. Wujud kapisan nyatakaké yèn
$\Gamma\Entails !A$ njalari $\Gamma\Proves !A$. Wujud kapindho
nyatakaké yèn $\Gamma$ konsistèn, $\Gamma$ bisa disukupi.
Kita mbuktèkaké wujud kapindho lan nyimpulaké wujud kapisan
saka asil iku. Buktiné mbutuhaké sawetara langkah. Wiwitané
ana himpunan konsistèn~$\Gamma$. Tambahna simbol konstanta
anyar $c_i$ cacah tanpa wates, lan formula awujud
$\lexists[x][!A(x)]\lif !A(c)$, kanthi masangaké saben
formula~$!A(x)$ kang mawa variabèl bebas ing basa kang wis
ditambahi karo salah siji konstanta anyar. Asilé himpunan ukara
\emph{jenuh} lan konsistèn kang ngemot~$\Gamma$.
Banjur ambaaké himpunan mau dadi \emph{himpunan jangkep lan
konsistèn}. Kanggo sembarang ukara~$!A$, himpunan jangkep lan
konsistèn ngemot salah siji saka $!A$ utawa $\lnot !A$.
Amarga wiwitané himpunan jenuh, saiki kita nduwèni himpunan
ukara~$\Gamma^*$ kang jenuh, jangkep, lan konsistèn, lan
ngemot~$\Gamma$. Saka himpunan iki bisa ditemtokaké struktur
$\Struct{M}$ kanthi $\Sat{M(\Gamma^*)}{!A}$ yèn lan mung yèn
$!A\in\Gamma^*$. Mligi, $\Sat{M(\Gamma^*)}{\Gamma}$, dadi
$\Gamma$ bisa disukupi. Yèn ana simbol $=$, konstruksiné rada
luwih ruwet.

Ana rong korolari wigati saka teorema kejangkeping.
\emph{Teorema kompaktitas} nyatakaké yèn $\Gamma\Entails !A$
yèn lan mung yèn $\Gamma_0\Entails !A$ kanggo sawijining
$\Gamma_0\subseteq\Gamma$ kang winates. Wujud ekuivalèné
nyatakaké yèn $\Gamma$ bisa disukupi yèn lan mung yèn saben
$\Gamma_0\subseteq\Gamma$ kang winates bisa disukupi.
Teorema iki migunani kanggo mbuktèkaké anané struktur kang
nduwèni sipat tartamtu. Upamané, saben teori kang nduwèni
modèl winates kanthi gedhé tanpa wates ndhuwur uga nduwèni
modèl tanpa wates. Dadi, kawinatesan ora bisa diandharaké
ing logika tataran kapisan. Korolari kapindho, \emph{teorema
L\"owenheim--Skolem}, nyatakaké yèn saben $\Gamma$ kang bisa
disukupi nduwèni modèl kang bisa dietung. Asil iku nuduhaké
yèn ora bisa dietung uga ora bisa diandharaké ing logika tataran
kapisan.
